% =====================================================================
% tesi/capitoli/20-caso-provenienza.tex
% =====================================================================
% copre: poke-ace/README.md
% copre: poke-ace/STUDIO-01-ace-e-legalita-in-home.md
% copre: poke-ace/STUDIO-02-marchi-di-origine-e-che-cosa-conta-una-collezione.md
% copre: poke-ace/STUDIO-03-la-risposta-della-comunita-e-le-due-severita.md
% copre: recreate-pokemon-distributions-events/STUDIO-04-campagna-di-trasferimento-e-il-tracciatore.md
% copre: recreate-pokemon-distributions-events/CATALOGO-EVENTI.md
% copre: pokedex-home-completo/README.md
% copre: pokedex-home-completo/STUDIO-01-che-cosa-vincola-la-scadenza.md
% copre: pokedex-home-completo/STUDIO-02-salvataggi-esterni-e-che-cosa-provano.md
% copre: pokedex-home-completo/CENSIMENTO-SALVATAGGI.md
% copre: pokedex-home-completo/STUDIO-03-la-catena-e-viva-e-la-lista-di-spunta.md
% copre: pokedex-home-completo/CHECKLIST-COMPLETA.md
% copre: generation-from-switch/README.md
% copre: generation-from-switch/STUDIO-01-scambio-automatico-e-provenienza.md
% verificato-al-commit: a427431
% ---------------------------------------------------------------------

\chapter{Caso di studio: la provenienza di un esemplare}
\label{cap:provenienza}

Il settimo caso di studio non ha per oggetto un apparecchio, un protocollo o un formato, ma una proprietà che nessuno dei capitoli precedenti ha dovuto nominare: la \emph{provenienza} di un dato, cioè non che cosa esso contenga ma per quale via sia arrivato dove si trova. Fino a qui il criterio di correttezza è stato interno al dato, e verificabile su di esso: un checksum torna o non torna, una struttura decifrata è coerente o non lo è, un valore individuale sta o non sta nell'intervallo ammesso. Questo capitolo tratta il caso in cui il criterio di correttezza è esterno al dato, e in cui due sequenze di byte identiche possono avere esiti opposti a seconda di come sono state prodotte.

La circostanza che rende il problema concreto è che l'obiettivo di lungo periodo enunciato nel capitolo~\ref{cap:distribuzioni}, cioè il completamento di una collezione dentro il deposito in rete della piattaforma corrente, ha una porzione che nessuna quantità di gioco può coprire: gli anni in cui non si è giocato, e gli eventi a tempo la cui finestra è chiusa da vent'anni. Per quella porzione le vie esistenti sono quattro, e differiscono fra loro non per efficacia ma appunto per provenienza.

\section{Le quattro vie, e ciò che le distingue}

La prima via è quella del capitolo~\ref{cap:distribuzioni}: si ricostruisce il meccanismo originale di distribuzione e si fa rifare al gioco, su hardware posseduto, esattamente ciò che il gioco faceva allora. L'esemplare che ne risulta è coerente per costruzione, poiché il procedimento che lo genera è quello che lo generava.

La seconda via scrive i byte dell'esemplare direttamente nel salvataggio, sfruttando l'esecuzione di codice arbitrario descritta nel capitolo~\ref{cap:esecuzione}. L'esemplare che ne risulta è coerente rispetto ai controlli che il suo costruttore conosceva, che è una proprietà più debole e la cui differenza dalla precedente il capitolo~\ref{cap:conversione} ha già enunciato.

La terza via non produce nulla in locale: riceve un esemplare da un terzo attraverso lo scambio, cioè attraverso il meccanismo che il gioco prevede per farlo. Chi la percorre non conosce la costruzione di ciò che riceve.

La quarta via è quella ufficiale, la catena dei trasferimenti fra generazioni, e produce esemplari autentici al prezzo di richiedere l'hardware intermedio di ogni anello e di dipendere da un servizio che ha una data di chiusura.

La distinzione che conta non è dunque fra vie lecite e illecite, che è una classificazione contrattuale, ma fra vie che producono un dato la cui storia coincide con la sua forma e vie che producono un dato la cui storia è diversa da ciò che la forma dichiara. Il resto del capitolo mostra che il servizio di destinazione, contro l'attesa ingenua, conosce la differenza.

\section{La domanda che decide, e la sua risposta in tre parti}

La domanda è se un esemplare la cui provenienza non sia una partita giocata venga accettato come legittimo dal deposito in rete, e se il tentativo esponga a una sanzione. La risposta si articola in tre parti, e nessuna delle tre è un sì o un no.

La prima parte è tecnica e favorevole alla fattibilità. I byte di un esemplare di terza generazione possono essere resi identici a quelli di un esemplare autentico, perché la struttura è documentata, la cifratura e il checksum sono calcolabili con il macchinario del capitolo~\ref{cap:cifratura}, e il metodo di generazione pseudocasuale è noto evento per evento, come il capitolo~\ref{cap:distribuzioni} ha stabilito. Un verificatore che guardi soltanto il dato non trova nulla da obiettare.

La seconda parte rovescia la prima, e la sua fonte è l'autore stesso degli strumenti che rendono possibile la via \cite{mankeymite-home}. Il deposito in rete conserva sul proprio lato l'informazione di quale via un esemplare abbia impiegato per entrarvi. Un esemplare autentico di terza generazione entra attraverso il servizio di deposito della console a doppio schermo; uno che venga dalla riedizione per la console corrente entra per una via diversa. Ne segue che, anche a parità di dati sottostanti, non è garantito che l'esemplare risultante sia indistinguibile da uno passato per l'altra porta. La coerenza del dato non implica la coerenza della storia, e il servizio conosce la storia.

Questa è la ragione per cui la provenienza è una grandezza e non una sfumatura morale. In tutti i capitoli precedenti l'informazione disponibile al verificatore era contenuta nel dato; qui il verificatore possiede un archivio proprio, alimentato a ogni transito, e quel canale è per costruzione inaccessibile a chi produce il dato. Nessun perfezionamento della costruzione può colmare quella asimmetria, perché non si tratta di un controllo più severo sul medesimo oggetto ma di un oggetto in più.

La terza parte è la politica dichiarata dal titolare del servizio, e va citata perché è l'unico elemento non congetturale sul rischio \cite{tpc-dati-alterati}. Chi risulti impiegare dati alterati per mezzi non autorizzati può subire la restrizione del gioco in rete, la restrizione delle funzioni di scambio nella versione per dispositivo mobile, e la sospensione dell'accesso al deposito nelle versioni per console e per dispositivo mobile, in forma temporanea o indefinita a discrezione del titolare e senza rimborso. La medesima fonte dichiara un'eccezione che delimita il rischio: non vi sono restrizioni per chi possieda dati alterati senza intenzione, per esempio ricevendoli in uno scambio senza saperlo.

Quella eccezione va letta per ciò che dice. Descrive chi riceve senza sapere, e chi si rivolge deliberatamente a un servizio di generazione sa esattamente che cosa sta ricevendo: non è dunque la persona di cui la clausola parla. Trattarla come un'assoluzione sarebbe la lettura interessata di un testo che è chiaro, ed è il genere di ragionamento contro cui la gerarchia delle fonti del capitolo~\ref{cap:fonti} è stata costruita.

Sulla verificabilità pratica la risposta è netta e va scritta come tale: la questione non è verificabile al momento della redazione. La compatibilità fra la riedizione dei due titoli della terza generazione e il deposito in rete non è ancora in funzione, e i controlli che il servizio applicherà a un esemplare proveniente da là non sono noti perché non esistono ancora. Non è nemmeno noto se il servizio accetterà esemplari di provenienza dichiarata Rubino, Zaffiro o Smeraldo che si trovino dentro un salvataggio della riedizione, oppure se rifiuterà le combinazioni che il titolare consideri impossibili.

\section{I marchi di origine, e la definizione di collezione completa}

La sezione precedente afferma che il servizio conosce la provenienza, sulla scorta della testimonianza di un autore competente. Resta da stabilire in quale forma la conosca, perché fra un'annotazione interna agli archivi e un contrassegno visibile a chiunque guardi l'esemplare la differenza è grande: la prima si potrebbe ignorare, il secondo no. La risposta si ottiene da uno strumento di comunità che modella quella informazione, il cui contenuto non è servito da alcuna interfaccia di dati ed è stato ricavato sondando per costanti il suo codice compilato \cite{home-checklist}. Vale ricordare il grado di fiducia che ne discende: quanto segue è il modello dei dati di una implementazione di terzo livello, non una dichiarazione del titolare del servizio.

Il modello nomina undici \term{marchi di origine}, che sono contrassegni categorici associati al gioco da cui un esemplare proviene: l'assenza di marchio, il marchio delle riedizioni virtuali della prima e della seconda generazione, quello della sesta generazione, quello dei titoli della settima, quello delle riedizioni della prima generazione per la console corrente, quelli dell'ottava generazione nelle sue tre diramazioni, quello della nona, quello del titolo più recente, e un marchio dedicato alla terza generazione. Il marchio del gioco per dispositivo mobile è trattato a parte nel codice.

Due dettagli di quel codice contano più dell'elenco. Il primo è che l'insieme completo dei marchi è definito una volta e, immediatamente accanto, ne è definito un secondo che è il primo meno il solo marchio della terza generazione: l'implementazione distingue dunque i marchi che esistono da quelli già ottenibili, e l'unico tenuto fuori è quello della porta che si apre a ottobre 2026. Il secondo è che, nella tabella che associa i giochi ai marchi e alle collezioni, alla riedizione della terza generazione corrisponde quel marchio e nessun altro.

La conseguenza per il caso di studio è netta e sfavorevole alla seconda e alla terza via. Un esemplare di terza generazione che raggiunga il deposito per la catena storica non porta alcun marchio, poiché quella catena precede l'introduzione dei marchi; uno che entri dalla porta nuova porterà il marchio dedicato. I due non saranno dunque indistinguibili nemmeno all'occhio di chi li guarda, e non soltanto negli archivi del servizio: la provenienza sarà stampata sull'esemplare. La strategia implicita in ogni costruzione a posteriori, cioè confidare che nessuno guardi, è resa impraticabile non da un controllo più severo ma dal fatto che l'informazione è pubblica.

Va detto con la medesima precisione ciò che questo non dimostra. Non dimostra che il servizio rifiuterà quegli esemplari, e non dimostra che il marchio sia il criterio con cui li valuterebbe, poiché un marchio dice da quale porta un esemplare è entrato e non se i suoi dati siano coerenti. Dimostra che la provenienza è osservabile, che è una proprietà più debole della prima e più forte di quanto la sezione precedente potesse affermare da sola.

Il medesimo sondaggio produce un risultato che riguarda l'obiettivo di collezione anziché la sua legittimità, e che va enunciato perché corregge un modo di porre la domanda. Quello strumento non definisce un unico insieme da completare: definisce quindici profili predefiniti più uno personalizzato, e sopra il profilo agiscono interruttori indipendenti che moltiplicano il conteggio. Il registro vivente ordinario conta una casella per specie; il registro per marchio di origine conta una casella per specie e per marchio, cioè fino a undici caselle per la medesima specie; un terzo interruttore ignora del tutto il marchio; le differenze di sesso si contano in due modalità; le forme alternative si contano o no, con due specie a cardinalità sproporzionata dotate di interruttore proprio. A queste dimensioni si aggiungono undici collezioni di provenienza speciale, fra cui quella delle distribuzioni di evento che è il bersaglio del capitolo~\ref{cap:distribuzioni} e quelle dei due titoli con esemplari oscuri, che nessuna delle quattro vie di questo capitolo copre. E si aggiungono quattro livelli di reperibilità, tutti attivi per difetto, di cui il quarto è dichiaratamente ipotetico: una collezione che li includa tutti non è completabile per definizione.

Ne segue che la domanda se una data lista sia la lista completa di ciò che si può possedere non ha risposta, non per insufficienza di dati ma per malposizione: fra il profilo minimo e quello per marchio di origine la cardinalità differisce di un ordine di grandezza, e nessuno dei due è più corretto dell'altro. Un obiettivo di collezione enunciato come superlativo non identifica dunque un insieme, e per essere verificabile deve nominare le proprie dimensioni: quali marchi, quali collezioni, quali forme, quali differenze di sesso e quali livelli di reperibilità. È una decisione e non un dato, e sta a chi conduce il lavoro prenderla.

Va infine registrato ciò che il sondaggio non ha letto, per non attribuirgli più di quanto abbia dato: il catalogo delle singole specie non compare in chiaro nel codice compilato e nessun servizio esterno lo fornisce, quindi questa sezione documenta la tassonomia e le regole, non l'elenco. Una domanda su una singola specie o forma non si risponde da qui, e la via più economica per ottenere l'elenco non è decodificare il codice ma esportare il foglio di calcolo che quello strumento stesso produce.

\section{La porta nuova, e la correzione del calendario}

Il capitolo~\ref{cap:distribuzioni} enuncia una scadenza esterna e la tratta come assoluta. Un annuncio del titolare della serie, del 13 agosto 2026, la rende assoluta soltanto in parte, e la correzione va fatta qui perché cambia che cosa è urgente \cite{frlg-home}.

L'annuncio dichiara che le riedizioni per la console corrente dei due titoli della terza generazione si collegheranno al deposito in rete a ottobre 2026, con la versione 4.1.0 del servizio, e porta tre dettagli che vale registrare: il trasferimento è a senso unico, cioè un esemplare che lascia quei titoli non può rientrarvi e nessun esemplare di altri titoli può visitarli; la capienza del piano a pagamento sale da seimila a novemila esemplari; e il completamento dell'inventario di quei titoli consegna un esemplare di una specie altrimenti non ottenibile.

Ne discende una revisione della pianificazione in tre punti. Il primo è che la chiusura del servizio di deposito della console a doppio schermo resta la scadenza ultima per la prima, la seconda, la quarta e la quinta generazione, la cui catena passa necessariamente da esso; per la sola terza generazione si apre una seconda porta che non ne dipende. Il secondo è che i due passaggi interni alla catena, dalla terza alla quarta generazione attraverso la struttura di migrazione e dalla quarta alla quinta attraverso il trasferimento senza fili, sono funzioni locali dei titoli e continueranno a operare dopo la chiusura: ciò che cessa è il solo tratto finale, e la corsa contro il tempo riguarda dunque quel tratto e non l'intera catena.

Il terzo punto è quello che si decide invece di constatarlo, ed è che la porta nuova non entra nella pianificazione come se funzionasse. Non si sa quali controlli il servizio le applicherà, la medesima fonte che ne annuncia l'apertura raccomanda di trasferire ora per la via ufficiale ciò che si può trasferire ora, e il trasferimento verso quei titoli è dichiarato a senso unico. Il costo della prudenza è nullo se la porta funziona, mentre il costo dell'assunzione opposta, se non funziona, è la perdita definitiva di tutto ciò che sarebbe passato per l'altra. Va registrato inoltre un vincolo che pesa sul tratto finale più di quanto sembri: il trasferimento dal servizio di deposito verso il deposito in rete richiede il piano a pagamento di quest'ultimo, mentre il servizio in chiusura è gratuito, e il piano gratuito del deposito conserva trenta esemplari. Una catena completata fino al servizio in chiusura non è dunque una collezione al sicuro.

\section{La primitiva della seconda via, e ciò che le è stato costruito sopra}

La tecnica su cui poggia la seconda via è quella del capitolo~\ref{cap:esecuzione}, dove serve a rimuovere un esemplare dal salvataggio di partenza; qui interessa la sola parte che è propria di questo uso.

Nella terza generazione la tecnica sfrutta punti in cui il gioco tratta come indirizzo o come istruzione un dato che l'utente controlla, e il vettore consueto sono i nomi delle scatole del deposito, che sono stringhe scritte dal giocatore e residenti in memoria in una posizione raggiungibile. La conseguenza che conta è che quelle stringhe diventano un canale di ingresso per codice: si scrive nei nomi delle scatole una sequenza che il gioco esegue, e quella sequenza scrive dove le si dice di scrivere.

Sopra quella primitiva la comunità ha costruito ciò che rende la tecnica utile anziché soltanto notevole, e la fonte lo chiama scrittore in base sessantaquattro: un programma che, una volta in esecuzione, accetta dati in una codifica digitabile e li scrive nel salvataggio byte per byte. Da quel momento il problema non è più far eseguire codice, ma comporre i byte giusti, e comporre i byte giusti è precisamente il problema che questo lavoro ha già risolto per altre ragioni, con la referenza dei formati e con il pacchetto software del capitolo~\ref{cap:ponte}.

\section{Gli strumenti della comunità, e una assenza significativa}

L'inventario che segue proviene dal canale dedicato di un server di community e riporta la funzione dichiarata di ciascuno strumento; nessuno di essi è stato eseguito nel corso di questo lavoro, e per la maggior parte le pagine non sono state aperte. La distinzione va mantenuta, secondo la gerarchia del capitolo~\ref{cap:fonti}: sapere che uno strumento esiste e che cosa dichiara di fare non è averlo verificato.

Sulla generazione del codice da eseguire esistono due generatori, uno generale e uno specifico per i due titoli della riedizione, che producono la sequenza da digitare a partire dall'effetto desiderato; sono il primo anello della catena e ciò che rende la tecnica accessibile senza scrivere assembly a mano. Sulla composizione dei dati esistono un convertitore da esadecimale alla codifica digitabile, che alimenta lo scrittore descritto sopra, e uno strumento che converte la richiesta di un oggetto in nomi di scatola, il quale è il caso semplice della medesima tecnica e conviene percorrere prima di applicarla a un esemplare. Il più rilevante per questo lavoro è però un costruttore di esemplari di terza generazione, che compone il dato completo comprese le vecchie distribuzioni di evento, con le opzioni di generazione conforme attive per difetto e con l'avvertenza, dichiarata dal suo autore, che l'accettazione da parte del deposito in rete non è garantita \cite{gen3-ace-builder}.

Sulle procedure esistono due guide di allestimento, una per le versioni non giapponesi della riedizione prima dei capi dei Quattro e una per gli altri due titoli della generazione, più un archivio che raccoglie codici cercabili per titolo, lingua e piattaforma, guide, domande frequenti e risoluzione dei problemi \cite{ace-archive}. Su quell'archivio va registrata un'assenza, perché è significativa quanto un contenuto: esso non contiene alcuna dichiarazione sulla legittimità degli esemplari prodotti, sulla loro accettazione da parte dei verificatori, né sui rischi per l'account. Tratta la tecnica come problema di implementazione e non di conseguenze, il che è una scelta legittima per un archivio tecnico e un difetto per chi vi cercasse la risposta della sezione precedente.

Una convergenza va enunciata perché cambia le opzioni di due capitoli. Il costruttore di esemplari genera anche le vecchie distribuzioni di evento, cioè il medesimo risultato che il capitolo~\ref{cap:distribuzioni} persegue ricostruendo la ROM di distribuzione e inviandola per multiboot. Le due vie producono lo stesso esemplare per strade opposte: una fa rifare al gioco ciò che il gioco faceva, l'altra scrive il risultato. La prima è più lenta e richiede hardware; la seconda non ha la storia. Quale convenga dipende interamente dalla risposta della sezione precedente, ed è un esempio limpido di come una domanda aperta governi una scelta di metodo invece di limitarsi a lasciarla incerta.

\section{La terza via: ricevere invece di produrre}

La terza via merita una sezione propria perché sposta il problema dal piano tecnico a quello della responsabilità, e la fonte che la riguarda è il titolare del servizio e non la comunità.

Nei titoli per la console corrente la comunità ha costruito bot di scambio: programmi che stanno in ascolto, ricevono una richiesta, generano l'esemplare richiesto e lo consegnano al richiedente attraverso il meccanismo di scambio del gioco \cite{berichandev}. Chi li ospita li fa girare su console modificata; chi li usa si limita a compiere uno scambio, ed è precisamente questa asimmetria a rendere la via interessante e insidiosa insieme.

L'utilità per l'obiettivo di collezione è precisa e limitata: quei servizi coprono le specie e le forme delle generazioni recenti, cioè quelle che un giocatore assente in quegli anni non possiede e per le quali non esiste alcun evento passato da ricostruire. Sono dunque complementari e non alternativi alle due vie che coprono le generazioni antiche, e non risolvono la questione della catena di trasferimento poiché operano già a valle di essa \cite{amiibodoctor-generazione}.

Sul piano dei dati vale quanto già stabilito: un esemplare ricevuto è coerente rispetto ai controlli che il suo costruttore conosceva, e chi lo riceve non sa nemmeno quali fossero. Sul piano della responsabilità la posizione differisce dalle altre vie, ma non nella direzione che tornerebbe comoda: l'eccezione della politica dichiarata copre chi riceve senza sapere, e chi si rivolge deliberatamente a un servizio di generazione non è in quella condizione. Il rischio residuo non è quindi nullo e resta dell'ordine di quello della seconda via, poiché la sanzione dichiarata riguarda l'accesso al servizio che custodisce la collezione.

Cinque verifiche restano da compiere, nessuna delle quali richiede hardware: se gli esemplari così ottenuti superino un verificatore di conformità indipendente; quali specie e forme quei servizi coprano effettivamente, che è la domanda che dice se la via serva all'obiettivo o soltanto lo tocchi; se l'esemplare ricevuto porti un marchio di origine e quale, poiché il marchio di origine è la parte visibile dell'informazione di provenienza \cite{home-checklist}; se il canale che ospita il servizio sia effettivamente quello indicato e in quali termini; e se la documentazione di terze parti disponibile sia ancora attuale, dato che un anno in questo dominio è molto. Fino a quelle verifiche la via ha una domanda e non una conoscenza, e presentarla altrimenti sarebbe riempire con ipotesi lo spazio che una lettura dovrà occupare.

\section{Una via composta, marcata come inferenza}

Questa sezione contiene un'inferenza di questo lavoro e non una citazione, ed è marcata come tale perché va verificata prima di essere impiegata per decidere; se regge, è la conclusione più utile del capitolo.

Il capitolo~\ref{cap:ldn-trading} descrive lo scambio fra un calcolatore e la riedizione dei due titoli attraverso il protocollo di rete locale, cioè un canale che fa entrare strutture di terza generazione nel titolo su console per il meccanismo di gioco previsto. Il capitolo~\ref{cap:smeraldo} stabilisce che quelle strutture si ottengono dal salvataggio di una cartuccia posseduta, letto con il lettore descritto là. Componendo i due fatti si ottiene una via che nessuna delle fonti nomina.

\begin{quote}
cartuccia posseduta $\rightarrow$ lettura del salvataggio $\rightarrow$ estrazione della struttura dell'esemplare $\rightarrow$ scambio via rete locale verso la riedizione $\rightarrow$ da ottobre 2026, trasferimento diretto nel deposito in rete
\end{quote}

Il pregio di quella via rispetto alla scrittura diretta dei byte è che il dato non è generato ma autentico: proviene da una cartuccia posseduta, con la sua storia vera, e il suo ingresso nel titolo su console avviene per scambio, cioè per il meccanismo previsto. Resta la medesima incertezza della sezione seconda su ciò che il servizio registrerà come provenienza, poiché anche qui l'ingresso avviene dalla porta nuova; ma la differenza fra un dato autentico consegnato per una via non ufficiale e un dato costruito è la stessa che questo lavoro documenta fra coerenza per costruzione e coerenza rispetto ai controlli noti.

Va dichiarato ciò che la rende non ancora praticabile, perché sono tre condizioni e nessuna è piccola: il lettore di cartucce non è ancora disponibile, il codice del capitolo~\ref{cap:ldn-trading} non è iniziato e la sua praticabilità dipende da una misura sull'adattatore senza fili che non è stata compiuta, e il collegamento con il deposito non esiste prima di ottobre. La via è dunque un'ipotesi di lavoro con tre precondizioni, non un piano.

\section{Le tre severità, e perché la domanda era mal posta}

La sezione sulla domanda che decide, all'inizio di questo capitolo, presuppone un verificatore soltanto e chiede se un esemplare costruito lo superi. Quella formulazione va corretta, e la correzione discende da una fonte di quinto livello interrogata dopo la stesura, cioè il complesso delle testimonianze della comunità che impiega la tecnica \cite{mankeymite-server}. I verificatori sono tre, con severità decrescente, e la domanda utile non è se un esemplare passi ma quale dei tre lo esamini.

Il più severo è il verificatore di conformità della comunità. Esige la coerenza fra il valore di personalità, l'identificativo dell'allenatore e quello segreto, e rileva un esemplare a cui la natura sia stata modificata dopo la generazione, poiché la modifica della natura altera il valore di personalità e rompe la corrispondenza con il seme e con il fotogramma dell'incontro.

Il secondo è il deposito in rete, la cui severità è descritta dalle fonti con una formula che vale riportare: rileva gli errori clamorosi, per esempio un luogo di incontro impossibile, e non l'incoerenza fine che il primo trova. Ne segue la conseguenza che appariva paradossale e non lo è: un esemplare può essere legale per il deposito e non legittimo per il verificatore della comunità, giacché le due domande sono diverse.

Il terzo, e in senso pratico il più severo, non è un programma ma il gioco competitivo ufficiale, dove un esemplare alterato comporta la squalifica.

La distinzione che ne discende è quella fra \term{legale} e \term{legittimo}, e non è una sfumatura terminologica: essa determina quale delle due proprietà si stia perseguendo, e un obiettivo di collezione che non la dichiari non è verificabile. Il presente lavoro adotta la seconda come propria, con la conseguenza che la conformità al deposito non è un criterio sufficiente.

\section{Il tracciatore, e la differenza fra un cancello e un'esposizione}

Alla proprietà osservabile descritta nella sezione sui marchi di origine si aggiunge un secondo meccanismo, che le fonti nominano separatamente e che ha conseguenze di natura diversa: a ogni esemplare che entra nel deposito viene assegnato un identificativo univoco di tracciamento. La conseguenza riferita è che un esemplare duplicato prima dell'ingresso non è distinguibile, mentre uno che entri, riceva il tracciatore e venga duplicato dopo risulta segnalato \cite{mankeymite-server}.

La rilevanza per un obiettivo di lungo periodo è maggiore di quella dei marchi, e conviene enunciarla nella forma più netta. Un marchio è un dato statico che dichiara la porta di ingresso; un tracciatore rende ogni esemplare depositato identificabile in modo persistente, il che significa che l'esame non avviene una volta sola al momento dell'ingresso ma resta possibile indefinitamente. Le fonti aggiungono, indipendentemente e a mesi di distanza, la clausola che completa il quadro: la valutazione corrente è che gli esemplari in questione siano legali per il deposito, e questo potrebbe cambiare in futuro.

Il medesimo meccanismo è descritto in modo più esteso da una seconda fonte indipendente, di autore diverso e anteriore di un anno, che non cita la prima \cite{amiibodoctor-generazione}. Su una materia che il titolare del servizio non documenta, la concordanza di due testimonianze indipendenti è il grado di fiducia massimo ottenibile, e resta il quinto livello. La seconda fonte aggiunge il fatto che rende il tracciatore non aggirabile in linea di principio e non soltanto in pratica: esso non è un campo del dato ma un riferimento a un archivio che risiede presso il servizio, quindi un identificativo scritto a mano con un editor viene rilevato come falso, poiché il servizio non lo ha mai emesso. La medesima fonte ne dichiara tre usi: distinguere gli esemplari individuali, rilevare i caricamenti duplicati provenienti da più account, e ricostruire lo stato di un esemplare gioco per gioco quando esso si sposta fra titoli.

Da queste affermazioni discende una distinzione che nessuna delle due fonti enuncia e che va enunciata qui, poiché ometterla lascerebbe a portata di mano la conclusione sbagliata. Il tracciatore non ostacola un esemplare che entri nel deposito attraverso una porta vera: chi ricostruisce un esemplare dentro la riedizione per console moderna e lo invia al deposito non falsifica alcun tracciatore, giacché il servizio glielo assegna in quel momento come lo assegnerebbe a qualunque altro. Ciò che il tracciatore impedisce è altro, cioè fabbricare un esemplare la cui provenienza dichiarata sia incompatibile con l'assenza di storia, e duplicare un esemplare dopo l'ingresso. La seconda fonte colloca in quella prima categoria gli esemplari dell'epoca della console a doppio schermo, per i quali non esiste alcuna strada che non passi dal deposito, e dichiara che la loro fabbricazione da zero tende a essere impedita allo scambio o segnalata presso il servizio.

Ne segue la formulazione corretta del rischio, che differisce da quella della sezione precedente e la sostituisce. Il rischio non è un cancello da superare una volta, ma un'esposizione permanente: ciò che oggi è accettato resta identificabile domani, e la regola con cui verrà esaminato domani non è quella con cui è stato accettato oggi. Per un obiettivo dichiarato come impresa di lungo periodo questa è la differenza fra un costo una volta pagato e un debito che resta aperto.

\section{Che cosa il deposito controlla, in cinque categorie}

La sezione sulle tre severità stabilisce che il deposito è il verificatore intermedio e ne caratterizza la severità con una formula, cioè che rileva gli errori clamorosi e non l'incoerenza fine. Quella formula si può sostituire con un elenco, che una fonte di quarto livello fornisce e che vale riportare perché due delle cinque categorie toccano campi che il presente lavoro compone con provenienza debole \cite{amiibodoctor-generazione}.

Va premessa la filosofia generale che la fonte enuncia, poiché è il criterio con cui interpretare l'elenco: l'idea è che un esemplare passi se avrebbe potuto concepibilmente essere ottenuto in modo legittimo nel gioco da cui dichiara di provenire. La fonte aggiunge da sé che esistono eccezioni e che il controllo non è perfetto.

La compatibilità di gioco è la prima e la più semplice: un esemplare si sposta soltanto in un titolo dove quella specie esiste, e un titolo incompatibile lo rifiuta.

La legalità del contenitore di cattura è la seconda, e porta un dettaglio di rilievo per questo lavoro. Il deposito conserva il contenitore originale e lo mostra correttamente al proprio interno; un gioco di destinazione che non lo preveda può mostrarne uno generico o uno dichiaratamente estraneo, e riportando l'esemplare nel deposito ricompare l'originale. La rilevanza è diretta, poiché il contenitore è il campo con la provenienza più debole fra quelli che lo strumento descritto nel capitolo~\ref{cap:strumenti} compone: sapere che il servizio lo conserva e lo espone significa che un valore errato non passa inosservato.

I dati di evento sono la terza, e la fonte li indica come la causa tipica del rifiuto di un esemplare composto male. Il deposito conosce quali distribuzioni di evento sono esistite e impone i vincoli evidenti: il contenitore di pregio, i formati fissi del nome e dell'identificativo dell'allenatore, il contrassegno dell'incontro fatidico, e certi riconoscimenti che devono o non devono essere presenti. È l'insieme dei campi che il capitolo~\ref{cap:distribuzioni} ricostruisce, e la coincidenza fra ciò che il servizio controlla e ciò che quel capitolo determina non è casuale: entrambi guardano ciò che l'evento fissava.

Le sequenze di mosse sono la quarta, e il meccanismo è più fine di quanto si supporrebbe: il deposito conserva le mosse separatamente per ciascun gioco, e quando un esemplare raggiunge un titolo dove le sue mosse correnti sono impossibili esse vengono sostituite con mosse valide, mentre il ritorno nel titolo precedente ripristina la sequenza che quello conosceva. Ne segue che l'insieme delle mosse non è una proprietà dell'esemplare ma una funzione della coppia formata dall'esemplare e dal titolo, il che è la seconda delle due categorie che toccano un campo di provenienza debole nel nostro strumento.

I blocchi sulla lucentezza e le restrizioni sugli esemplari mitici sono la quinta: il deposito e il sistema di scambio globale filtrano le combinazioni possibili nel codice ma non realmente disponibili, per esempio le distribuzioni a lucentezza bloccata. La fonte osserva che questa categoria appare gestita manualmente dal titolare e non è sempre corretta, il che è a sua volta informativo: significa che l'insieme delle combinazioni rifiutate non è derivabile da una regola e va trattato come un elenco che può contenere errori in entrambe le direzioni.

\section{La via composta, non più un'inferenza}

La sezione dedicata alla via composta la presentava come inferenza di questo lavoro non sostenuta da alcuna fonte, e la marca va rimossa. La medesima via è descritta dall'autore degli strumenti della comunità, in una forma più semplice di quella qui ipotizzata: per la terza generazione si ricostruisce una copia esatta dell'esemplare all'interno della riedizione per console moderna mediante l'esecuzione di codice, dopo averne letto il valore di personalità e gli altri dati dalla cartuccia, e da là si trasferisce al deposito \cite{mankeymite-server}.

Le due semplificazioni rispetto alla forma ipotizzata sono sostanziali. Non occorre il protocollo di rete locale, poiché l'esemplare non viene trasferito nella riedizione ma ricostruito dentro di essa; e il presupposto non è possedere un canale di comunicazione ma possedere i dati, cioè avere letto la cartuccia originale. La via del protocollo di rete resta disponibile come alternativa, ed è nominata dalla medesima fonte insieme al progetto che la implementa, con l'osservazione che per quella via non serve alcun codice.

Che tale via sia praticabile dipende da un fatto verificato in questa medesima fonte: l'esecuzione di codice funziona sulla riedizione per console moderna. All'annuncio della riedizione la comunità dubitava che vi funzionasse; nei mesi successivi sono comparse guide e testimonianze d'uso. Due particolari operativi vanno registrati perché non sono deducibili. Il primo è che i codici differiscono da quelli per cartuccia, poiché la riedizione applica ai nomi delle scatole un filtro sulle parole vietate che sposta le sequenze digitabili. Il secondo è che quel filtro non è noto con esattezza nemmeno all'autore degli strumenti, il quale dichiara che l'elenco impiegato dal proprio costruttore è una ricostruzione empirica accumulata sui casi incontrati e non corrisponde a quello del titolare: è una ammissione che conferma per via indipendente ciò che il confronto sul codice di quello strumento aveva mostrato a proposito della sua tabella di codifica dei caratteri.

\section{Un campo che non si falsifica per natura, e non per difficoltà}

Le sezioni precedenti riportano che il tracciatore assegnato dal deposito non è falsificabile e che la sua assenza è rilevabile. Resta da enunciare perché, poiché la ragione non è di grado ma di specie e distingue quel campo da tutti gli altri trattati in questo lavoro.

Ogni altro campo di un esemplare è calcolabile, e la parte tecnica di questo lavoro consiste nel saperlo calcolare. Il checksum è una somma di parole che si ricalcola; il valore di personalità è una parola che si compone da due estrazioni; i valori individuali si impaccano cinque bit per volta; natura, sesso, indice di abilità e lucentezza discendono per formula dal valore di personalità, come l'appendice~\ref{app:generatori} svolge per esteso. Per ciascuno di essi esiste una regola, e chi la conosce la soddisfa: la conferma è che un'implementazione indipendente, ricevuto un esemplare così composto, ne ha ricostruito il seme.

Il tracciatore non ha una regola da soddisfare. Non è una funzione dei dati dell'esemplare, e chi lo legge non lo verifica ricalcolandolo: è una chiave che rimanda a un archivio detenuto dal servizio, il quale la valuta cercandola e confrontando ciò che vi trova con l'esemplare che ha davanti. Un valore inventato non è dunque un valore sbagliato ma un valore inesistente, e viene rilevato come tale non perché sia mal formato ma perché il servizio non lo ha mai emesso.

La distinzione ha una conseguenza che vale enunciare in forma generale, poiché si applica oltre il caso. Un campo la cui correttezza è definita da un calcolo è falsificabile da chiunque conosca il calcolo, e la sua difesa dipende dalla segretezza della regola, che è una difesa fragile e che questo lavoro documenta essere stata superata su ogni altro campo. Un campo la cui correttezza è definita da una corrispondenza con uno stato remoto non è falsificabile da nessuno che non abbia accesso a quello stato, e la sua difesa non dipende da alcun segreto: nessun progresso nell'analisi del formato la scalfisce, perché non c'è nulla da analizzare. Il primo tipo di campo protegge finché la regola è ignota; il secondo protegge per costruzione.

Il verso opposto completa il quadro e discende dalla medesima natura: un campo di quel tipo non si può nemmeno lasciare in bianco dove il contesto lo attende, poiché la sua assenza è già una risposta. Ne segue che, per un esemplare la cui provenienza dichiarata implichi il transito dal deposito, non esistono tre possibilità ma una: o vi è transitato, e allora il campo è pieno di un valore emesso; o non vi è transitato, e allora nessuna scrittura locale rimedia.

Vale osservare, per non trarne una conclusione più larga del dovuto, che tutto ciò riguarda un esemplare che dichiari una provenienza incompatibile con l'assenza di storia. Un esemplare che entri nel deposito attraverso una delle porte effettive riceve il proprio tracciatore dal servizio nel momento dell'ingresso, e per esso la questione non si pone: la sezione precedente distingue le due situazioni, e confonderle porta a credere chiuse anche le vie che sono aperte.

\section{Le tre destinazioni possibili, e perché una di esse non esiste}

Le sezioni precedenti stabiliscono che il servizio conserva un tracciatore per ogni esemplare e che esso non è falsificabile. Resta da trarne la conseguenza operativa, ed è più netta di quanto la formulazione consueta della domanda suggerisca. La domanda consueta è quale scadenza si applichi a un esemplare che si voglia portare nel deposito in rete; la risposta è che la scadenza dipende dalla porta, e che una delle tre destinazioni che sembrano disponibili non è una porta.

\subsection*{Scrivere l'esemplare in un salvataggio della generazione corrente}

Fra un titolo della terza generazione e uno della nona non esiste alcun collegamento di gioco: la catena che li congiunge passa necessariamente dal deposito. Scrivere direttamente l'esemplare dentro un salvataggio della generazione corrente è dunque un'operazione sui byte e non un trasferimento, e produce un esemplare che si trova dove non poteva arrivare.

Che ciò sia rilevabile è stato dimostrato per via strumentale nel capitolo~\ref{cap:distribuzioni}: il verificatore, leggendo un esemplare di terza generazione nel contesto di un titolo della nona, ha dichiarato mancante il tracciatore. Lo cerca perché in quel contesto è atteso, e l'esemplare non lo porta perché non è transitato da nulla.

Ne segue una chiusura in entrambi i sensi, ed è la ragione per cui questa destinazione non è una porta. Con il tracciatore non si può, poiché esso non è un campo del dato ma un riferimento a un archivio che risiede presso il servizio, e un identificativo scritto a mano è rilevato come falso perché il servizio non lo ha emesso. Senza il tracciatore non si può, poiché la sua assenza costituisce essa stessa un'obiezione. Non è una via che scade: è una via che non esiste, e la scadenza del servizio di deposito non ha alcun rapporto con essa.

\subsection*{Le due porte che esistono, e le loro scadenze disuguali}

L'esemplare va collocato dove appartiene, ossia in un salvataggio della propria generazione, e da là le porte sono due con scadenze diverse.

La prima è la catena storica descritta nel capitolo~\ref{cap:distribuzioni}. Di essa cessa il solo tratto conclusivo, alla data dichiarata; i due passaggi interni fra terza, quarta e quinta generazione sono funzioni locali dei titoli e continuano a operare, quindi la scadenza si applica all'ultimo tratto e non all'intera catena.

La seconda è il collegamento diretto fra la riedizione dei due titoli della terza generazione e il deposito, che non dipende dal servizio in chiusura e per il quale non è dichiarata alcuna scadenza. Per impiegarlo l'esemplare deve trovarsi in un salvataggio di quei due titoli, e le due strade per portarvelo sono quelle discusse in questo capitolo, ossia ricostruirvelo mediante l'esecuzione di codice oppure trasferirvelo attraverso il protocollo di rete locale. In entrambi i casi l'ingresso nel deposito avviene da una porta effettiva e il tracciatore viene assegnato dal servizio, il che rimuove l'obiezione della sottosezione precedente per costruzione anziché per aggiramento.

\subsection*{Il campo di cui la scelta dipende, e una mossa che rimuove un'incognita}

Il rapporto del verificatore contiene, fra le righe informative, il titolo di origine dichiarato dall'esemplare. Nel caso esaminato esso era il titolo che il corpus di metadati propone come valore predefinito, e va osservato che il medesimo corpus non dichiara per quell'evento alcun insieme di titoli ammessi: il valore non è dunque vincolato dalla fonte.

La scelta non è però indifferente rispetto alla seconda porta, poiché quella porta appartiene ai soli due titoli della riedizione. Un esemplare che dichiari un titolo diverso dovrebbe comunque trovarsi in un salvataggio di quei due, e su quella condizione precisa le fonti registrano un punto aperto: non è noto se il servizio accetti un esemplare la cui provenienza dichiarata sia un titolo diverso da quello della riedizione in cui si trova.

Ne discende una mossa che elimina l'incognita anziché aggirarla, e che vale enunciare come criterio generale: dove un campo non è vincolato dalla fonte e da esso dipende quale via l'esito potrà impiegare, quel campo va scelto in coerenza con la via che si intende impiegare, e non lasciato al valore predefinito. Dichiarare come titolo di origine uno dei due della riedizione non è una furbizia sul dato, poiché entrambi i valori sono ammessi dalla fonte: è la scelta coerente fra valori equivalenti per la fonte e non equivalenti per la destinazione.

\section{La dimensione del problema, e la necessità di una selezione}

Un dato pratico chiude il capitolo e riguarda l'obiettivo anziché la tecnica. L'autore degli strumenti dichiara in quattrocento il numero degli esemplari da distribuzione di cui tenere traccia, e aggiunge il commento che vale come pianificazione: conviene scegliersene alcuni, poiché altrimenti non si termina prima della chiusura del servizio di deposito \cite{mankeymite-server}.

Quel numero è una dichiarazione di quinto livello e non è stato verificato; la conseguenza che ne discende non dipende però dalla sua esattezza. L'insieme completo delle distribuzioni non è raggiungibile nel tempo residuo, quindi una selezione è necessaria, e l'unica scelta effettiva riguarda se compierla deliberatamente o subirla quando il tempo finisce. È l'elemento che mancava alla decisione sul profilo di collezione discussa nella sezione sui marchi di origine, e la rende più concreta: non si tratta soltanto di dichiarare quali dimensioni contare, ma di ordinare per priorità all'interno di una dimensione che non si può esaurire.

Va infine registrata un'attesa della comunità, etichettata come tale perché non è stata verificata su fonte ufficiale: si attende una riedizione per console moderna dei titoli di Hoenn, che renderebbe legittimi per quella via anche gli esemplari di quella regione. Il presente lavoro non la adotta come pianificazione e la registra come dato sul contesto. Nella medesima fonte compare tuttavia un'avvertenza che conviene raccogliere, poiché proviene da chi conosce gli strumenti meglio di chiunque: la domanda se convenga astenersi dall'impiegare il codice che completa l'inventario nazionale finché quella riedizione non esista, per possibili conseguenze sul deposito. È la medesima prudenza che questo lavoro ha adottato riguardo alla porta di ottobre, formulata da chi ha più ragioni di volere il contrario.

\section{La strategia, e i punti che restano aperti}

La conclusione operativa non è di questo lavoro ma della fonte che ha aperto la questione, e coincide con ciò che il capitolo~\ref{cap:distribuzioni} aveva stabilito per altra via, il che è la migliore conferma disponibile.

Ciò che si può trasferire ora per la via ufficiale si trasferisce ora, e non si attende alcun aggiramento per gli esemplari che hanno una strada aperta, perché quella strada ha una data di scadenza mentre gli aggiramenti hanno un rischio non quantificato. Ciò che nessuna via legittima può dare è il solo caso in cui la seconda e la terza via hanno senso: gli anni non giocati e gli eventi passati sono quel caso. La porta nuova si prova quando esiste, su materiale sacrificabile e non sulla collezione, verificando che cosa il servizio accetti prima di esporvi qualunque cosa a cui si tenga. E la decisione sull'impiego delle vie non ufficiali si prende esplicitamente, perché il rischio dichiarato non è la perdita di un esemplare ma la sospensione dell'accesso al servizio che custodisce la collezione: è una sola decisione per entrambe le vie, poiché espongono il medesimo account alla medesima sanzione, e deciderle separatamente produrrebbe la contraddizione di accettare il rischio per una e rifiutarlo per l'altra a parità di esposizione.

Cinque punti restano aperti, e l'ordine in cui sono elencati è quello del loro potere di decidere. Se il deposito in rete accetti esemplari prodotti per queste vie e con quali controlli: non verificabile prima di ottobre 2026, ed è la domanda che decide tutto il resto. Se accetti esemplari di provenienza dichiarata Rubino, Zaffiro o Smeraldo che si trovino in un salvataggio della riedizione: la fonte lo dichiara ignoto. Che cosa esattamente il servizio conservi come informazione di provenienza e se sia consultabile dall'utente: la fonte afferma che esiste ma non ne descrive la forma, e i marchi di origine ne sono la parte visibile. Se la via composta della sezione precedente regga, con le sue tre precondizioni. E se il costruttore di esemplari produca, per le distribuzioni di evento, dati coincidenti con quelli che il capitolo~\ref{cap:distribuzioni} ricostruisce dal metodo di generazione: è un confronto fattibile senza hardware e senza esporre alcun account, e falsificherebbe o confermerebbe entrambe le vie in un colpo, il che lo rende il passo successivo di maggior valore fra quelli disponibili.

Vale chiudere osservando che questo capitolo non risolve il proprio problema, e che il modo in cui non lo risolve è esso stesso un risultato. La domanda della sezione seconda ha un'istruttoria chiusa e una risposta impossibile: tutto ciò che si può sapere è scritto, e ciò che manca non è conoscenza ma un evento futuro. Distinguere le due condizioni è ciò che permette di non ripetere l'indagine credendo di poterla chiudere, ed è la ragione per cui la sezione dedicata ai punti aperti dichiara, per ognuno, non soltanto che è aperto ma perché.
\section{La provenienza storica, che è una categoria di fatto a sé}

Il capitolo ha trattato fin qui la provenienza come proprietà interna di un esemplare, cioè come l'insieme dei campi che dichiarano da quale gioco e da quale incontro esso derivi. Esiste una seconda accezione della medesima parola, che il presente lavoro ha incontrato tardi e che merita una sezione perché appartiene a una categoria epistemica differente: la provenienza come fatto storico, cioè il luogo, la data e il modo in cui un esemplare venne effettivamente consegnato a chi lo ricevette.

\subsection*{Perché nessuna fonte di primo livello la contiene}

La gerarchia delle fonti adottata da questo lavoro colloca al primo livello i disassemblati e le implementazioni di riferimento, cioè il codice, e ai livelli inferiori la documentazione derivata. Per i fatti meccanici la gerarchia è una guida sicura, poiché un disassemblato è il gioco stesso e non può sbagliarsi su ciò che il gioco fa. Per i fatti storici la medesima gerarchia è inutilizzabile, e la ragione è di principio: nessun disassemblato può sapere in quali negozi un dono fu distribuito, né in quali giorni, poiché quelle informazioni non sono nel programma. Un codice non contiene la storia della propria diffusione.

Ne segue che per questa categoria di fatti la fonte migliore disponibile è di livello inferiore, tipicamente un'enciclopedia collaborativa, e che ciò non costituisce un ripiego ma il riconoscimento di un limite della gerarchia. La conseguenza operativa adottata è tenere le due categorie in file distinti: i fatti meccanici si rigenerano dalla fonte di primo livello a ogni corsa, i fatti storici sono autorati con il collegamento e la data di lettura accanto a ciascuna voce, e dove le due fonti divergono la divergenza viene scritta anziché risolta in silenzio, con la regola che sui fatti meccanici prevale il codice e sui fatti storici prevale l'unica fonte che se ne occupi.

Una convalida incrociata merita di essere registrata, poiché è il tipo di riscontro che vale più di entrambe le fonti prese singolarmente. Le pagine enciclopediche lette dichiarano, per ciascuna distribuzione, l'identificativo numerico dell'allenatore; la tabella dell'implementazione di riferimento dichiara il medesimo dato, e le due sono state ricostruite da persone diverse a partire da materiali diversi. Su trentasette gruppi esse concordano su tutti gli identificativi confrontabili tranne uno. La discordanza residua è stata risolta con un argomento interno anziché per autorità: gli identificativi di questo catalogo codificano la data nella forma anno, mese e giorno, e il valore dichiarato dal codice si legge come una data coerente con la manifestazione a cui la distribuzione è attribuita, mentre l'altro no.

\subsection*{La durata della finestra, e perché è la misura della rarità}

Fra i fatti che la provenienza storica porta e che nessuna ispezione del dato può suggerire, il più utile alla pianificazione è la durata della finestra di distribuzione. Essa varia, nel catalogo esaminato, di quattro ordini di grandezza: da tre ore in un solo giorno per una distribuzione tenuta in una catena di negozi statunitensi, a tre anni e due mesi per la campagna che correggeva un difetto dell'orologio interno delle prime cartucce.

La conseguenza è quantitativa e non retorica. Detta $\lambda$ la frequenza media con cui gli esemplari venivano consegnati e $T$ la durata della finestra, il numero di esemplari autentici esistenti è dell'ordine di $\lambda T$, e poiché $\lambda$ varia entro limiti stretti, essendo governata dalla capienza di un banco e dalla pazienza di una coda, la variabilità di $T$ domina interamente il risultato. Ne segue che la reperibilità odierna di un esemplare autentico è approssimativamente proporzionale alla durata della sua finestra, e che conoscere quella durata dice, prima di ogni ricerca, quali esemplari convenga cercare e quali convenga ricreare. È un caso in cui un fatto storico, che non ha alcuna influenza sui byte, determina interamente la strategia con cui i byte verranno prodotti.

\subsection*{Accettato e fedele, che non sono la stessa cosa}

L'ultima distinzione è quella che questa sezione esiste per enunciare, e il presente lavoro l'ha formulata soltanto dopo averne violato il lato più difficile.

Un esemplare può essere valutato secondo due criteri indipendenti. Il primo è la legittimità, cioè la proprietà di essere accettato da un verificatore di conformità: essa riguarda i soli campi che il verificatore vincola, ed è una proprietà decidibile, perché il verificatore la decide. Il secondo è la fedeltà, cioè la proprietà di coincidere con l'esemplare che fu effettivamente distribuito: essa riguarda tutti i campi, compresi quelli che nessun verificatore controlla, e non è decidibile da alcun programma, poiché richiede di sapere che cosa esistette.

I due criteri divergono su un insieme di campi che è facile individuare in astratto e facile dimenticare in pratica: quelli che il gioco lascia modificare dopo la consegna. L'oggetto tenuto è l'esempio canonico. Un verificatore non può pretendere che un esemplare porti l'oggetto con cui fu distribuito, poiché un oggetto si toglie, si scambia e si consuma, e un esemplare privo di esso è indistinguibile da uno che l'abbia semplicemente usato; ne segue che il campo non è vincolato e che un esemplare senza oggetto passa i controlli. Il caso concreto incontrato è che la distribuzione italiana della campagna del decennale consegnava il proprio Pikachu con un oggetto specifico, documentato dalla sola fonte enciclopedica: la tabella dell'implementazione di riferimento non lo dichiara, il verificatore non lo pretende, e gli esemplari prodotti da questo lavoro ne erano privi senza che nulla lo segnalasse.

Vale enunciare il criterio generale che ne discende, poiché si applica a qualunque ricostruzione e non a questa sola. Quando si misura la qualità di una ricostruzione con il giudizio di un verificatore, si sta misurando la legittimità e non la fedeltà, e la differenza fra le due è precisamente l'insieme dei campi che il verificatore ignora. Quell'insieme non è vuoto e non è piccolo, e i suoi elementi hanno una proprietà comune che li rende insidiosi: sono ignorati dal verificatore proprio perché sono modificabili, e sono modificabili proprio perché appartengono alla storia dell'esemplare anziché alla sua identità. Chi voglia la fedeltà deve quindi documentare la storia, e la documentazione della storia non può venire dal codice.
\section{Una scadenza che non vincola ciò che sembrava vincolare}

Il capitolo ha distinto la legittimità dalla fedeltà, e la specie dall'esemplare. Questa sezione mostra che cosa quella distinzione produca quando la si applica a un vincolo di tempo, e il risultato è che il vincolo si sposta interamente da una delle due categorie all'altra.

\subsection*{Il problema, e la lettura ingenua}

Un servizio di deposito accetta esemplari per due vie. La prima proviene dai titoli recenti e non ha termine. La seconda passa per un servizio in dismissione, con una data di cessazione annunciata. La lettura ingenua del vincolo è che la data limiti la collezione, e da essa segue il piano di lavoro ovvio, cioè spendere il tempo residuo sulla collezione.

La lettura è sbagliata, e la struttura dell'errore è la seguente: essa confonde l'insieme degli oggetti da collezionare con l'insieme degli oggetti che soltanto la via in dismissione può fornire. Il secondo è la differenza fra il primo e l'unione delle vie permanenti, e la sua cardinalità non è nota a priori. Può essere l'intero insieme, può essere vuota, e da quale dei due casi si dia dipende interamente il piano.

\subsection*{Il calcolo, e il dato su cui va fatto}

La cardinalità si calcola. Detto $U$ l'insieme delle voci collezionabili, $D_i$ l'insieme delle voci fornite dal titolo permanente $i$ e $I_j$ quello fornito dal titolo in dismissione $j$, l'insieme vincolato dalla scadenza è
\begin{equation}
V = \Bigl(\bigcup_j I_j\Bigr) \setminus \Bigl(\bigcup_i D_i\Bigr) ,
\end{equation}
e la domanda pratica è la cardinalità di $V$ e non quella di $U$.

Il calcolo richiede, per ciascun titolo, l'elenco delle voci che esso contiene. Su questo il presente lavoro ha commesso un errore di ricognizione che vale raccontare, poiché è di una specie ricorrente. La prima ispezione ha cercato quel dato fra i file che l'implementazione di riferimento dedica agli incontri, essendo l'incontro il modo in cui una voce si ottiene; li ha trovati in numero di duecento, in un formato binario compresso con struttura diversa per generazione; e ne ha concluso che il dato non fosse ottenibile a costo ragionevole, dichiarando formalmente una deroga alla regola del lavoro che vuole i dati generati e non trascritti.

La conclusione era falsa perché la domanda era mal posta. Non serviva sapere quali incontri un titolo offra, ma quali voci esso contenga, che è un dato diverso e sta altrove: nelle tabelle delle statistiche di base, che sono array di record a dimensione fissa con un contrassegno di presenza in un bit noto, e si leggono in cinquanta righe senza riscrivere alcun lettore. La deroga è stata ritirata poche ore dopo essere stata dichiarata.

La lezione generale è che la difficoltà percepita di un dato dipende da quale domanda si è formulata, e che davanti a un dato apparentemente inaccessibile conviene chiedersi se esista una domanda più debole che serva allo stesso scopo. Nel caso in esame la domanda debole, cioè la presenza anziché l'ottenibilità, era sufficiente per il calcolo di $V$ e costava due ordini di grandezza in meno.

\subsection*{Il risultato, e la correzione che lo completa}

Il calcolo dà $|V| = 0$ al livello delle specie: l'unione dei titoli permanenti copre l'intero insieme delle milleventicinque voci del catalogo nazionale. Al livello delle forme la differenza contiene dodici elementi, dei quali dieci sono forme totemiche di una generazione e due forme di sola battaglia.

Le seconde sono escluse per definizione, poiché una forma che esiste soltanto durante un incontro non può stare in un deposito. Le prime richiedono una correzione a una conclusione che la ricerca preliminare aveva tratto, ed è la parte più istruttiva del caso. Quella ricerca le dichiarava irrecuperabili dopo la scadenza, il che implica che fossero recuperabili prima. Il verificatore dei trasferimenti dell'implementazione di riferimento afferma invece due cose: che quattro di quelle specie non si trasferiscono affatto alla generazione successiva, e che per le altre l'esemplare trasferito deve portare la forma base anziché quella totemica, pena l'invalidità.

Ne segue che nessuna di quelle voci può esistere nel deposito, né prima né dopo la data. Esse non sono in attesa di scadenza: sono irraggiungibili per costruzione, e la loro assenza dalla collezione non è un costo della dismissione ma una proprietà permanente del formato. Il risultato completo è dunque $|V| = 0$ a entrambi i livelli.

Vale isolare la forma logica della correzione, poiché non riguarda un conteggio ma una modalità. Dire che un oggetto è irrecuperabile dopo una data significa affermare implicitamente che sia recuperabile prima, e quella affermazione implicita è una proposizione distinta che richiede una verifica propria. Il presente lavoro ha già incontrato tre volte l'errore di ricavare un'informazione positiva da un'assenza; questo ne è una quarta variante, in cui da una scadenza si ricava una disponibilità.

\subsection*{Dove il vincolo si è spostato}

L'insieme $V$ è vuoto, ma il vincolo di tempo non è scomparso: si è spostato interamente sull'altra categoria, cioè sui singoli esemplari la cui identità richiede una provenienza anteriore. Essi non appartengono a $U$ e non si contano in voci di catalogo, poiché non sono specie ma esemplari: la specie è sostituibile e l'esemplare no.

La conclusione operativa è dunque l'inversione della lettura ingenua. Il tempo residuo non va speso sulla collezione, che non ha scadenza, ma sugli esemplari, che ne hanno una e per i quali nessuna via alternativa esiste. È esattamente il lavoro che i capitoli precedenti descrivono, e la sezione dà a esso, retrospettivamente, la sua giustificazione quantitativa: quel lavoro non era una parte del problema fra le altre, era l'unica parte con una scadenza.

\subsection*{Una nota sul valore di una fonte, che dipende dalla domanda}

Il conteggio è stato confrontato con una base dati indipendente, quella che serve un simulatore di battaglie della comunità. Il confronto ha prodotto un accordo esatto sul numero delle specie, portando a tre le ricostruzioni indipendenti concordi, e una discrepanza sul numero delle forme.

La discrepanza non è un errore di nessuna delle due parti, e la sua spiegazione è un'osservazione generale sull'uso delle fonti. Quella base dati serve il gioco competitivo, dunque enumera le forme che modificano il comportamento di un esemplare in battaglia e omette quelle puramente estetiche; le tabelle dell'implementazione di riferimento enumerano invece tutto ciò che esiste nei dati. Le due contano insiemi diversi, e chi le confrontasse senza saperlo attribuirebbe a un errore ciò che è una differenza di definizione.

Ne segue che una fonte non è semplicemente attendibile o inattendibile: è attendibile su certe domande e muta su altre, e la sua utilità dipende dal riconoscere quali. Impiegare quella base dati per contare le forme non produce un errore evidente ma un numero preciso, ricavato correttamente, e non pertinente, che è il genere di risultato più difficile da individuare in una revisione.
\section{Un artefatto di provenienza esterna, e le tre domande che pone}

Il capitolo si chiude con una classe di materiale che i precedenti non avevano incontrato e che merita una trattazione propria, perché il modo tipico di sbagliare su di essa è un errore di categoria e non di calcolo. Il materiale è un salvataggio prodotto da qualcun altro e reso disponibile in rete; l'errore è trattare come una domanda sola tre domande distinte.

La prima domanda è se il file sia strutturalmente un salvataggio funzionante del gioco che dichiara di essere. Ha una risposta binaria, poggia su somme di controllo, e si risponde con un programma.

La seconda è che cosa contenga, cioè quali esemplari e dunque quali voci del catalogo contribuirebbe. Ha una risposta numerica, poggia su un censimento, e si risponde con lo stesso programma.

La terza è se sia lecito impiegarlo. Non ha una risposta che un programma possa dare, perché non riguarda i byte del file ma le conseguenze del suo uso su un account e su una console, e la sua sede è un documento di norme. Uno strumento che rispondesse anche a questa nasconderebbe una decisione dentro un calcolo, ed è la ragione per cui lo strumento costruito qui chiude ogni corsa dichiarando esplicitamente ciò che il proprio esito non dice.

\subsection*{Il rischio effettivo, e perché la rassicurazione consueta è la sbagliata}

La reazione consueta davanti a un file scaricato è passarlo a un antivirus, e l'esito pulito viene letto come esclusione del rischio. Vale mostrare perché in questo caso quell'esito, benché vero, sia informativo quasi quanto nulla.

Un salvataggio è dato e non codice. Nessun processo lo esegue: viene letto da un gioco che ne interpreta i campi a offset noti, oppure da un editor. Non può quindi ospitare un programma nel senso in cui lo ospita un eseguibile, e uno scanner non trova nulla perché in un file di quella natura non c'è nulla di quel genere da trovare. L'esito conferma una proprietà che era vera per costruzione, prima di qualunque scansione.

I rischi effettivi sono tre e nessuno è visibile a uno scanner. Il primo è la malformazione: un file che il gioco rifiuta all'avvio, o che accetta e poi corrompe. Si esclude verificando le somme di controllo che il formato porta, e questa è la ragione per cui la verifica strutturale, che sembra un dettaglio di ingegneria, è invece l'unico controllo di sicurezza che qui abbia senso. Il secondo è l'esemplare costruito per innescare un difetto del gioco, di cui la struttura corrotta che i giochi di terza generazione marcano con un bit dedicato è il caso classico: si riconosce dai campi incoerenti, e va detto con precisione che controllarli è meno di un giudizio di legittimità, non un suo equivalente. Il terzo non è tecnico, ed è il maggiore dei tre: impiegare esemplari altrui su un account proprio, che è la terza domanda e non la prima.

\subsection*{Identificare senza credere al nome}

Il nome di un file scaricato è una dichiarazione del suo autore, non un dato, e il metodo di identificazione impiegato qui è quello già usato in questo lavoro per stabilire il gioco di un salvataggio: la dimensione restringe la famiglia, e dentro la famiglia si prova ciascun candidato con il predicato che la fonte usa per distinguerli, riferendo quale risponde e quali no.

Vale l'osservazione generale che rende il metodo non negoziabile. Un predicato inventato che per caso funzioni sul file che si ha in mano è indistinguibile da uno corretto, e la differenza compare soltanto sul file successivo: la trascrizione dei predicati dalla fonte non è dunque pedanteria bibliografica ma la sola difesa contro una identificazione che sembra funzionare.

I predicati impiegati mostrano una varietà istruttiva. Per la quarta generazione si legge il piede del blocco generale e si verifica che il campo della dimensione coincida con la dimensione del blocco e che la parola seguente sia una delle due costanti note; e poiché la dimensione di quel blocco differisce fra i tre gruppi di titoli, il medesimo predicato che accerta la validità identifica anche il gioco, che è un caso raro e comodo. Per la quinta è un codice ciclico a sedici bit calcolato sul piede, con due lunghezze diverse per i due gruppi. Per la sesta e la settima è una firma di coda comune, e a distinguere i titoli resta la sola dimensione. Per la prima generazione sono due liste a offset noti, di cui si verifica che il contatore stia entro il massimo e che il byte successivo sia il terminatore, con offset diversi fra edizioni internazionali e giapponesi: due controlli banali, e in congiunzione bastano.

Un caso di ambiguità reale vale segnalarlo perché mostra come si risolve. Due formati appartenenti a piattaforme diverse hanno esattamente la medesima dimensione, e la dimensione da sola non li separa. Si separano perché portano strutture diverse in posizioni diverse: l'uno una firma di coda, l'altro una somma di controllo e un'intestazione di scheda di memoria. Verificare entrambe costa due confronti e rimuove l'ambiguità; deciderla sulla dimensione l'avrebbe nascosta.

\section{Una normalizzazione che cancella ciò che distingue}

Il censimento condotto con lo strumento appena descritto ha rivelato un difetto non nei file esaminati ma nel codice che li esaminava, e il difetto merita una sezione perché è un esemplare puro di una classe di errore che questo lavoro ha incontrato più volte in forme diverse.

\subsection*{Il caso}

Due fonti indipendenti numerano le specie in modo diverso, e per metterle in corrispondenza si confrontano i nomi. Poiché due fonti scrivono lo stesso nome con convenzioni tipografiche diverse, il confronto avviene su una forma normalizzata, ottenuta portando il nome a minuscole e rimuovendo ogni carattere che non sia lettera o cifra.

Fra i caratteri così rimossi ci sono i due segni di sesso, e due specie della prima generazione differiscono per il solo segno finale. Le due normalizzazioni coincidono, la corrispondenza associa il primo dei due nomi e scarta il secondo, e il numero della seconda specie finisce a puntare sull'identificatore interno della prima.

La gravità va misurata per quello che è. Non è una imprecisione: è attribuire a un numero di catalogo la specie sbagliata, che su questo dato è il difetto peggiore concepibile, poiché ogni campo derivato successivo, cioè gruppo di crescita, coppia di abilità, punti potenza e tabella delle mosse, viene letto dalla riga di un'altra specie. E il commento della funzione che costruisce la corrispondenza avvertiva esplicitamente di questo pericolo, citando come esempio una coppia di specie su cui le due numerazioni divergono: la funzione portava dentro il difetto dal quale dichiarava di proteggere.

\subsection*{Perché è rimasto invisibile}

Le ragioni sono tre e si sommano.

La prima è che la riga colpevole non appartiene al dominio del problema. Il problema difficile era la corrispondenza fra due numerazioni, ed era stato studiato, documentato e verificato su due specie di controllo; la normalizzazione del testo era un dettaglio di implementazione, del genere che si scrive senza pensarci. È la stessa asimmetria descritta nel capitolo precedente a proposito di dove i difetti si concentrino, e qui si presenta in forma ancora più netta, perché lo stesso file conteneva il presidio e la falla.

La seconda è che la corrispondenza non era mai stata interrogata sulla riga difettosa. Nessuna voce del catalogo degli eventi riguarda quelle due specie, quindi il generatore, che è l'unico consumatore fino a quel momento, non aveva mai chiesto quella riga. Il difetto era latente, e il primo uso nuovo della funzione lo ha attivato: è la ragione per cui un difetto in una struttura condivisa non va classificato per il danno che ha fatto ma per il danno che può fare.

La terza è la forma in cui si è manifestato, che è quella insidiosa. Non un errore, non un'eccezione, non un valore assurdo: un rapporto in cui un conteggio valeva trecentottantacinque anziché trecentottantasei, con una riga che segnalava una specie priva di corrispondenza. Un numero leggermente sbagliato dentro un output corretto è precisamente ciò che supera una revisione a video, e in questo caso è stato notato solo perché la specie mancante era implausibile: una specie comunissima, assente da una collezione che dichiarava di essere completa. L'implausibilità del risultato ha fatto da campanello, e va registrato che è un campanello che non suona sempre.

\subsection*{La correzione, e la parte che conta}

La correzione ha due parti, e la seconda vale più della prima.

La prima è tradurre i due segni in lettere distinte invece di cancellarli, così che le chiavi restino due. Risolve il caso.

La seconda è imporre e verificare le due proprietà che la corrispondenza deve possedere e che nessuno controllava. La prima proprietà è che la normalizzazione sia iniettiva sull'insieme dei nomi: due nomi distinti non devono produrre la medesima chiave. La seconda è che la corrispondenza stessa sia iniettiva: due numeri di catalogo non devono puntare al medesimo identificatore interno. Entrambe si esprimono in due righe, entrambe fanno terminare il programma con la ragione scritta, ed entrambe avrebbero colto questo difetto nel momento in cui fu introdotto anziché mesi dopo.

Il punto generale, che è la ragione per cui questa sezione esiste, è che le due proprietà erano note e implicite. Chiunque avesse riflettuto sulla struttura del problema le avrebbe enunciate senza esitazione, e nessuno le aveva scritte, perché sembravano troppo evidenti per meritare codice. Un invariante evidente e non verificato ha esattamente la stessa forza di un invariante falso, e la differenza fra i due si scopre soltanto quando qualcosa lo viola.

Una precisazione tecnica va aggiunta perché mostra che anche il controllo va calibrato. Applicato senza cautela, il controllo di iniettività della normalizzazione fallisce subito su due voci della tabella delle specie che non sono specie ma segnaposto per identificatori interni non impiegati, e portano come nome una sequenza di punti interrogativi: la loro normalizzazione è la stringa vuota, e la stringa vuota collide con se stessa. Vanno escluse prima del controllo, e la lezione minore è che un presidio troppo severo che vada rilassato dopo la prima corsa è comunque preferibile a un presidio assente, perché il rilassamento avviene con il caso concreto sotto gli occhi.

\section{Stato e esemplare, cioè due cose che il perimetro non distingueva}

Un'ultima distinzione chiude il capitolo, e nasce da un problema pratico: il primo anello della catena di trasferimenti richiede che una funzione del gioco sia stata sbloccata, e sbloccarla richiede di avere completato la storia principale di quel titolo. Con più titoli da attraversare, il costo di questa condizione domina il costo di tutto il resto.

La via che si offre è importare in quella cartuccia un salvataggio in cui la funzione risulti già disponibile. La norma di perimetro che regola i salvataggi di terze parti li ammette subordinando ogni esemplare importato al giudizio di un verificatore, esemplare per esemplare, e quella norma parlava di esemplari perché il caso che l'aveva generata riguardava esemplari.

Il caso presente non muove alcun esemplare. Ciò che entra è uno stato di avanzamento, cioè un insieme di indicatori che dichiarano un evento della storia come accaduto. Ne segue che il rischio che la norma intendeva prevenire, cioè l'impiego di esemplari di provenienza altrui, non si presenta: non perché sia stato ridotto, ma perché non ha oggetto.

La distinzione va enunciata come categoria e non come eccezione a un caso, perché ricorrerà. Sbloccare una funzione, possedere un oggetto necessario a un evento, avere ottenuto un riconoscimento sono tutti stati; nessuno di essi porta con sé la questione della provenienza di un'entità collezionabile. E la distinzione porta con sé le proprie condizioni, che sono ciò che la rende una distinzione onesta e non una scappatoia: lo spazio di deposito del salvataggio importato va svuotato prima dell'impiego, perché gli esemplari che contiene non servono allo scopo e la loro sola presenza reintrodurrebbe il problema; e le garanzie della norma sull'hardware restano intere, poiché l'operazione non è la lettura di un file ma la sovrascrittura di un supporto fisico.

Una nota conclusiva sull'asse che questa distinzione illumina. Il capitolo ha già separato legale da legittimo; il materiale di provenienza esterna introduce un terzo asse, indipendente da entrambi, che si può chiamare proprio. Un esemplare può essere legittimo, nel senso che nessun verificatore ha da contestargli, e non proprio, nel senso che porta il nome e l'identificatore di un altro e nel gioco figura come ricevuto. Un obiettivo di collezione formulato in termini di storia possibile riguarda il terzo asse almeno quanto il secondo, e chi non lo enuncia rischia di soddisfare una condizione credendo di soddisfarne un'altra.
\section{Contare non è censire, e perché la differenza emerge tardi}

Il capitolo si chiude con una osservazione di metodo che il lavoro ha imposto tardi e che sarebbe stata invisibile prima, poiché dipende dal numero delle fonti e non dalla loro qualità.

\subsection*{Il difetto di un progetto che conta}

Per settimane il lavoro qui descritto ha prodotto conteggi. Quante voci contenga un catalogo di distribuzioni; quante di esse un programma sappia riprodurre; quante specie stiano nei depositi di una raccolta di salvataggi; quante voci di dono portino le generazioni che dipendono da un servizio intermedio; quante specie e quante voci di forma le tabelle dei titoli dichiarino. Ogni conteggio era esatto, ciascuno era stato verificato contro la propria fonte, e nessuno rispondeva alla domanda che governa una campagna di collezione, cioè che cosa manchi.

La ragione è strutturale e vale enunciarla in forma generale: un conteggio è un numero e non un insieme, e i numeri non si intersecano. Sapere che una fonte offre $385$ elementi e un'altra $59$ non permette di dire quanti ne offrano insieme, poiché la loro sovrapposizione non è nel dato; e non permette di dire quanti restino, poiché non esiste il riferimento rispetto a cui calcolare il complemento. Due grandezze che sembrano confrontabili non lo sono, e la somma che verrebbe naturale scrivere è priva di significato.

Il censimento è l'operazione che rimedia, e richiede un oggetto che i conteggi non producono: un elenco esplicito delle voci desiderate, ciascuna con un'identità stabile, su cui ogni fonte si possa proiettare. Costruito quell'elenco, ogni conteggio precedente cessa di essere un numero isolato e diventa una colonna, cioè un predicato definito su un insieme comune; e la domanda su ciò che manca diventa una differenza insiemistica, che è calcolabile.

\subsection*{Perché la necessità emerge tardi, e non è una svista}

Vale difendere il lavoro precedente, perché la sua forma era adeguata al proprio momento. Con una o due fonti, ciascuna si descrive da sé e la sovrapposizione è ispezionabile a occhio; il costo di costruire un elenco di riferimento supererebbe il suo beneficio. La necessità compare quando le fonti diventano abbastanza numerose perché le loro intersezioni siano la parte interessante dell'informazione, e nel caso in esame ciò è avvenuto al passaggio da due a quattro famiglie di fonti.

Ne discende un criterio trasferibile, che è la ragione per cui questa sezione esiste. Il momento in cui passare dal conteggio al censimento non è dettato dalla dimensione del problema ma dal numero delle fonti che vi concorrono: finché sono poche, contare informa; appena si sovrappongono, contare smette di informare e continua a sembrare informativo. Il secondo è lo stato pericoloso, poiché produce rapporti pieni di numeri esatti e inutili.

\section{Un identificatore per una voce da ottenere}

L'elenco di riferimento richiede un identificatore, e la scelta apparentemente ovvia è sbagliata in un modo che vale analizzare, poiché l'errore è di tipo e non di merito.

\subsection*{Perché il numero di catalogo non serve}

Il catalogo nazionale assegna un numero a ogni specie, ed è l'identificatore che ogni fonte usa. Ha però una proprietà che lo rende inadatto allo scopo: è invariante rispetto alle distinzioni che qui contano. Non cambia per il sesso, non cambia per una variante regionale, non cambia per una forma alternativa. È esattamente ciò che si vuole da un identificatore di specie, ed esattamente ciò che non si vuole da un identificatore di voce da ottenere.

La conseguenza pratica è che una lista indicizzata per numero di catalogo non è una lista di ciò che si desidera ottenere ma di ciò che si desidera possedere almeno una volta, in qualche variante. Al termine di una campagna essa non distingue una voce soddisfatta da una soddisfatta in parte, e la distinzione è precisamente ciò per cui la lista esiste.

\subsection*{Le tre proprietà richieste}

L'identificatore adottato è la coppia ordinata fra numero di catalogo e indice di forma, resa come stringa a lunghezza fissa con la forma base a zero. Tre proprietà lo giustificano, e ciascuna esclude un'alternativa.

La prima è la stabilità rispetto alle implementazioni. Ogni programma che manipoli questi dati porta una propria numerazione interna, e questo lavoro ha già pagato il prezzo di confondere due numerazioni: un identificatore che poggiasse su una di esse erediterebbe quel rischio e lo propagherebbe a ogni documento derivato. La coppia poggia sulla sola numerazione pubblicata, che nessuna implementazione può cambiare.

La seconda è l'ordinabilità. La resa a lunghezza fissa con riempimento di zeri fa coincidere l'ordine lessicografico con quello del catalogo, cosicché ogni elenco ordinato per identificatore risulti già nell'ordine in cui un lettore lo cerca. È una proprietà minore che si paga con nulla e la cui assenza costa un riordinamento a ogni consultazione.

La terza è la totalità, ed è la più importante. L'identificatore esiste per ogni voce, comprese quelle di cui non si conosce il nome della forma; e poiché il nome della forma richiede una tabella di corrispondenza che varia per titolo e per lingua, un identificatore che lo contenesse sarebbe definito solo su una parte dell'insieme. Un'identità parziale non è un'identità.

\subsection*{Il limite dell'elenco, dichiarato dall'elenco stesso}

L'elenco enumera le voci di forma che i dati dichiarano, e per la maggior parte di esse non afferma se il deposito di destinazione le conti come caselle separate: nessuna documentazione di primo livello lo stabilisce. La scelta adottata è enumerare e marcare l'indeterminatezza anziché decidere, e vale enunciarne la ragione in forma generale, poiché la tentazione opposta è forte.

Un elenco di riferimento che risolvesse per proprio conto una questione aperta produrrebbe un documento internamente coerente e falso in un punto che nessuna verifica successiva potrebbe individuare, poiché la decisione sarebbe indistinguibile da un dato. Marcare l'indeterminatezza costa una colonna e conserva la distinzione fra ciò che si è letto e ciò che si è supposto, che è la sola proprietà che rende un documento correggibile.

\section{Una terza occorrenza, e il criterio che la governa}

Un difetto introdotto durante la costruzione di quell'elenco chiude il capitolo, perché è la terza istanza della medesima forma logica in tre giorni e la ripetizione permette di enunciare il criterio che la previene.

Il generatore dell'elenco necessitava di due insiemi che un altro programma già calcolava, e ne ha riscritto la lettura anziché invocarlo. Nella riscrittura, il caso particolare di un titolo la cui tabella non porta il contrassegno di presenza è stato risolto sostituendo l'insieme delle sue voci con l'insieme delle sue specie in forma base: una semplificazione che perde le voci di forma di quel titolo.

L'esito non è stato un errore ma un numero plausibile, inferiore di sei unità al corretto, con quattro voci che risultavano vincolate da una scadenza da cui non sono vincolate. La forma logica è la stessa dei due casi precedenti descritti in questo lavoro, cioè il bit di un'abilità dedotto per la via breve e una normalizzazione di testo che cancellava un carattere distintivo: una scorciatoia in una parte considerata ovvia produce un risultato che nessun controllo esistente distingue da quello giusto.

Il criterio che le governa tutte e tre non è una raccomandazione di attenzione, che non ha mai prevenuto nulla, ma una regola strutturale: quando due programmi hanno bisogno della medesima lettura, la lettura appartiene a uno solo dei due e l'altro la invoca. La correzione applicata non è stata correggere la riscrittura ma eliminarla, spostando la lettura in una funzione unica dove il trattamento del caso particolare vive accanto alla propria motivazione. E il difetto è stato individuato per una ragione che merita di essere registrata: i due programmi sono stati confrontati sui rispettivi numeri prima che al secondo si prestasse fede, e il confronto è costato una singola esecuzione.
